\documentclass[12pt]{article}
\usepackage[margin=1in]{geometry}
\usepackage[T1]{fontenc}
\usepackage[utf8]{inputenc}
\usepackage{lmodern}
\usepackage{microtype}
\usepackage{amsmath,amssymb,mathtools}
\usepackage{enumitem}
\usepackage{setspace}
\usepackage{xurl}
\usepackage[hidelinks]{hyperref}
\setstretch{1.60}
\setlength{\parindent}{0pt}
\setlength{\parskip}{6pt plus 2pt minus 1pt}
\setcounter{secnumdepth}{-\maxdimen}
\emergencystretch=3em
\sloppy
\begin{document}

\section{41. The Principal Genus Theorem for Relatively Galois Number Fields}

\textit{Original publication: Math. Ann. 108 (1933), pp. 411--419.}

Recent work has shown that non-commutative methods, and especially the theory of algebras, make it possible to formulate and prove known theorems on relatively cyclic and relatively abelian number fields for arbitrary relatively Galois fields. The chief example is the norm theorem, which in its general form is a theorem on split algebras. I shall show here that the principal genus theorem can be formulated in the same hypercomplex language: it is a theorem on inner automorphisms, or equivalently on crossed representations.

The proof is obtained from the theorem on split algebras by purely algebraic and arithmetic considerations. The transcendental core of that theorem is already the norm theorem in the cyclic case, and in fact it is enough there to treat prime degree. To make the formulation intelligible I first put before it a purely algebraic analogue, not used later in the proof, which I call the principal genus theorem in the minimal sense. In the cyclic special case it becomes the familiar theorem 90 of Hilbert's Zahlbericht: the norm of an element is one if and only if the element is of the form
\[
  a=b^{1-S}=b/b^S .
\]
I shall state this already as a theorem on inner automorphisms and on crossed representations, and prove it from the general facts on inner automorphisms and crossed product representations of algebras.

In the actual principal genus theorem the crossed product of a Galois group with the multiplicative group of a field is replaced by the crossed product of the ideal-class group and the Galois group. The factor systems are then divided into a finer induced ideal-class classification. This is the analogue of the ray-class division in class-field theory. Since the corresponding general theory of automorphisms and representations is not yet available in this form, the proof below reduces the statement to the theorem on split algebras. More concretely, it reduces the theorem for ideal classes to the corresponding theorem for ideals. At the ideal level Brandt's theorems on maximal orders in an algebra may be regarded as the substitute for the required automorphism theorems. They yield a proof parallel in spirit to the minimal proof, although complicated by the need to pass to the separate places. I give instead a nearly trivial proof communicated by Artin, which is an even simpler analogue of Speiser's argument.

\subsection{The principal genus theorem in the minimal sense}

Let \(k\) be an arbitrary commutative field, not necessarily a number field. Let \(K/k\) be a separable Galois extension of degree \(n\), and let \(G\) denote its Galois group. Recall first the standard facts on crossed products.

A crossed product is obtained by embedding simultaneously the field \(K\) and the group \(G\) into an algebra \(A\) in such a way that the automorphisms of \(K\) become inner automorphisms. Let
\[
  G=\{S,T,R,\ldots\}.
\]
Choose symbols \(u_S\), one for each element \(S\) of \(G\), and put first
\[
  A=\sum_{S\in G} u_S K
\]
as a right \(K\)-linear module. The requirement that \(u_S\) induce the automorphism \(S\) is expressed by
\[
  z u_S = u_S z^S\qquad(z\in K).
\]
Multiplication among the new basis symbols is prescribed by
\[
  u_S u_T = u_{ST} a_{S,T},\qquad a_{S,T}\in K^{\times} .
\]
Associativity is equivalent to the usual cocycle relation
\[
  a_{S,T}^{R} a_{ST,R}=a_{T,R}a_{S,TR}.
\]
With multiplication extended bilinearly, \(A\) is the crossed product
\[
  A=(a_{S,T},K,G).
\]
It is a central simple algebra over \(k\); \(K\) is a maximal commutative subfield and hence a splitting field. Conversely, for a given division algebra one obtains matrix algebras over it which can be represented in this way as crossed products.

Replacing \(u_S\) by
\[
  v_S=u_S c_S,\qquad c_S\in K^{\times},
\]
changes the factor system into the associated factor system
\[
  a'_{S,T}=a_{S,T}c_S^T c_T c_{ST}^{-1}.
\]
In particular, factor systems of the form
\[
  a_{S,T}=c_S^T c_T c_{ST}^{-1}
\]
are associated to the unit factor system. The quantities \(c_S^T c_T c_{ST}^{-1}\) are called transformation quantities. Associated factor systems are collected into a class; similarly, similar algebras are collected into a Brauer class. The two class notions correspond bijectively. With fixed splitting field \(K\), the classes form an abelian group under direct product; this is the Brauer group of algebra classes, and it corresponds to componentwise multiplication of cocycle classes. Its identity is the class of split algebras, or equivalently the class of all transformation systems.

For later comparison note the cyclic special case. Suppose \(K/k\) is cyclic and \(S\) is a generator of the Galois group. One can choose one symbol \(u\), with its powers corresponding to the powers of \(S\), and the crossed product is described by
\[
  A=K+uK+\cdots+u^{n-1}K,
\]
\[
  zu=uz^S,
\]
\[
  u^n=a,
\]
where \(a\in k^{\times}\). If \(v=uc\), then
\[
  a'=a\,N_{K/k}(c).
\]
Thus the group of cyclic algebra classes is represented by the factor group
\[
  k^{\times}/N_{K/k}(K^{\times}).
\]

The formulation of the minimal principal genus theorem uses the fact that the defining relations are purely multiplicative. They therefore also define a group extension
\[
  G^*=\{u_S K^{\times}: S\in G\}
\]
of \(G\) by \(K^{\times}\). The group \(K^{\times}\) is an abelian normal subgroup of \(G^*\), and
\[
  G^*/K^{\times}\cong G.
\]
The elements of \(G^*\) are exactly the regular elements of the algebra which transform \(K\) into itself. In other words they are exactly those \(g^*\) for which
\[
  (g^*)^{-1}Kg^*=K.
\]
They generate precisely the ring automorphisms of \(K\). Indeed, every such element gives an automorphism of \(K\), and conversely any such automorphism is represented by some \(u_S w\). If the induced automorphism on \(K\) is trivial, then \(w\) commutes with all elements of \(K\). Since \(K\) is a maximal commutative subfield, \(w\in K\), and because the element is regular, \(w\in K^{\times}\).

The minimal principal genus theorem can now be stated in three equivalent forms.

\textbf{First form.} Every automorphism of the group \(G^*\) which extends the identity automorphism of \(K^{\times}\) is inner, and is induced by an element of \(K^{\times}\).

\textbf{Second form.} If all transformation quantities
\[
  c_S^T c_T c_{ST}^{-1}
\]
have the value one, then the \(c_S\) are symbolic \((1-S)\)-st powers. Thus there exists \(b\in K^{\times}\) such that
\[
  c_S=b^{1-S}=b/b^S
\]
for every \(S\in G\).

\textbf{Third form.} The group \(G\) has only one crossed representation class of degree one in \(K^{\times}\) belonging to the unit factor system.

Here a representation \(u_S\mapsto c_S\) is called crossed with factor system \(a_{S,T}\) if
\[
  c_S^T c_T=c_{ST}a_{S,T}.
\]
Two such representations belong to the same class if
\[
  c_S=b^{-S}d_S b
\]
for an element \(b\in K^{\times}\). In degree one this is simply the usual equivalence by multiplication with a coboundary.

It remains to prove the first form and to see that the other two follow. Let an automorphism of \(G^*\), fixing \(K^{\times}\) elementwise, send \(u_S\) to \(v_S\). Since all the relations defining the crossed product are preserved, the direct correspondence
\[
  \sum_S u_S b_S\longmapsto \sum_S v_S b_S
\]
defines a ring automorphism of \(A\) which fixes every element of \(K\). Such an automorphism of a central simple algebra is inner. By the characterization of the normalizer of \(K\) just recalled, the implementing element must lie in \(K^{\times}\). This proves the first form.

For the second form, the relations force each \(v_S\) to have the shape \(u_S c_S\). Preservation of the factor system is exactly the assertion that the transformation quantities are one. Conversely, if this condition holds, then \(v_S=u_Sc_S\) defines an automorphism of \(G^*\) of the kind just considered. The first form gives
\[
  v_S=b^{-1}u_Sb=u_S(b/b^S),
\]
so \(c_S=b^{1-S}\). In the cyclic case the condition becomes \(N(c)=1\), and the conclusion \(c=b^{1-S}\) is Hilbert's theorem 90. The third form is merely the same statement expressed in representation language: a degree-one crossed representation with unit factor system is equivalent to the identity representation.

The third form is only a special case of a broader statement. Given any factor system \(a_{S,T}\), there is only one irreducible crossed representation class attached to it. Its degree is the Schur index of the corresponding algebra, that is, the degree of the associated division algebra. The present degree-one case is the part of this general assertion needed here.

\subsection{The principal genus theorem}

We now pass from the minimal theorem to the arithmetic theorem. In the cyclic form of the principal genus theorem of class-field theory two classifications occur at once: absolute ideal classes in the upper field and ray classes in the lower field. In the general Galois case this is replaced by the following mechanism. A classification of the ideals of the upper field induces a finer classification of the factor systems.

Let \(\mathcal I\) denote the group of all ideals of \(K\). Since \(G\) acts on ideals, one can form the same extension by the symbolic elements \(u_S\). Thus one has complexes \(u_S\mathcal I\) and multiplication rules parallel to the crossed-product rules. Passing from ideals to absolute ideal classes, write \(I\) for the ideal-class group. The complexes \(u_S I\) are still meaningful, since the principal class is stable under the action of \(G\). Hence the equivalence relation which passes from ideals to classes transfers from \(\{u_S\mathcal I\}\) to \(\{u_S I\}\). If \(\mathfrak a\) and \(\mathfrak b\) are equivalent ideals, then \(u_S\mathfrak a\) and \(u_S\mathfrak b\) are equivalent as well.

Consequently the symbols \(u_S\) are equivalent to all products \(u_S(c_S)\), where \((c_S)\) ranges over principal ideals of \(K\). If one requires the product relations to be single-valued in the induced class sense, one is requiring that the principal-ideal transformation systems
\[
  (c_S)^T(c_T)(c_{ST})^{-1}
\]
belong to the identity factor-system class. The identity class for factor systems is not merely the class of trivial systems; it consists of those principal-ideal factor systems \((a_{S,T})\) which contain representatives such that the associated algebras
\[
  (a_{S,T},K,G)
\]
split at all ramified places of \(K/k\). These principal ideals form a group under multiplication of factor systems. This group contains all transformation systems arising from principal ideals, because those give split algebras everywhere.

The principal genus theorem can now be stated exactly parallel to the minimal theorem.

\textbf{First form.} Suppose the induced ideal-class classification of factor systems is used. If the substitutions
\[
  v_S=u_S c_S
\]
define an automorphism of the complexes \(u_S I\), in the sense that the relations among the \(v_S\) agree with the original relations under the induced class division, then this automorphism is inner. It is generated by a single ideal class \(\mathfrak b\).

\textbf{Second form.} If the transformation classes
\[
  c_S^T c_T c_{ST}^{-1}
\]
all lie in the identity class of the induced factor-system classification, then the vector of ideal classes \(\{c_S\}\) belongs to the principal genus and is a symbolic \((1-S)\)-st power. Thus there exists an ideal class \(\mathfrak b\) such that
\[
  c_S=\mathfrak b^{1-S}=\mathfrak b/\mathfrak b^S
\]
for every \(S\in G\).

\textbf{Third form.} Relative to the induced ideal-class division of the factor systems, the group \(G\) has only one crossed representation class of degree one in the ideal-class group belonging to the identity class of factor systems.

The equivalence of these three formulations is the same as in the minimal case. The passage from the first to the second is even simpler here, because the automorphism has already been assumed to be generated by substitutions of the form \(v_S=u_Sc_S\). The assumption is necessary: the ideal-class group need not be a maximal commutative subgroup; for instance all classes may be ambiguous. In the third formulation the representation matrices are necessarily monomial matrices, since only multiplication is defined in the class group. For representations of degree one this makes no difference.

\subsection{Proof}

I prove the second form. The proof rests on two lemmas.

\textbf{Lemma 1 (the ideal form of the principal genus theorem).} If the ideals \(\mathfrak c_S\) have transformation systems equal to the unit ideal,
\[
  \mathfrak c_S^T\mathfrak c_T\mathfrak c_{ST}^{-1}=\mathcal O_K,
\]
then the \(\mathfrak c_S\) are symbolic \((1-S)\)-st powers. There is an ideal \(\mathfrak b\) such that
\[
  \mathfrak c_S=\mathfrak b^{1-S}
\]
for every \(S\in G\).

Equivalently, the first and third formulations hold at the ideal level, with the same proviso as above that automorphisms are represented by substitutions \(v_S=u_S\mathfrak c_S\).

The proof is immediate. Put
\[
  \mathfrak b=\sum_S \mathfrak c_S,
\]
where the sum is the module sum of ideals, equivalently the greatest common divisor in the sense of ideal theory. The displayed transformation equations imply
\[
  \mathfrak b^T=\sum_S \mathfrak c_S^T=\sum_S \mathfrak c_{ST}\mathfrak c_T^{-1}=\mathfrak b\mathfrak c_T^{-1}.
\]
Hence \(\mathfrak c_T=\mathfrak b/\mathfrak b^T=\mathfrak b^{1-T}\).

\textbf{Lemma 2 (the split-algebra theorem in the required form).} Let
\[
  A=(a_{S,T},K,G)
\]
be a crossed-product algebra. Suppose \(A\) splits at all finite and infinite ramified places of \(K/k\), and suppose the principal ideals \((a_{S,T})\) are transformation quantities of ideals:
\[
  (a_{S,T})=\mathfrak c_S^T\mathfrak c_T\mathfrak c_{ST}^{-1}.
\]
Then \(A\) splits everywhere, hence splits absolutely. Consequently the ideals \((a_{S,T})\) are transformation quantities of principal ideals:
\[
  (a_{S,T})=(d_S)^T(d_T)(d_{ST})^{-1}.
\]
Conversely, every everywhere split algebra satisfies the stated local condition.

This is just the familiar behavior of crossed products under passage to local fields. Let \(k_{\mathfrak p}\) be the completion of \(k\), and let \(K_{\mathfrak P}\) be the corresponding completion of \(K\) at a prime \(\mathfrak P\) over \(\mathfrak p\). Let \(Z\) be the decomposition group of \(\mathfrak P\), and let \(a_Z\) be the part of the factor system indexed by elements of \(Z\). Then the local algebra is
\[
  A_{\mathfrak p}\sim (a_Z,K_{\mathfrak P}/k_{\mathfrak p},Z).
\]
If \(\mathfrak p\) is unramified, the decomposition group is cyclic and the local algebra is the distinguished cyclic presentation attached to the unramified inertia field. Under the hypothesis of the lemma, the restricted factor system is associated to one consisting of local units. Passing to the normalized cyclic presentation, the parameter is a unit. In an unramified local extension all units that occur here are norms; hence the local cyclic algebra splits. At the infinite places splitting is part of the hypothesis. Thus \(A\) splits at every place. The global theorem on split algebras now gives global splitting.

The final step is immediate. Choose ideals \(\mathfrak c_S\) representing the ideal classes \(c_S\). The assumption that the class transformation systems belong to the identity class of the induced factor-system division says precisely that the ideal transformation systems
\[
  \mathfrak c_S^T\mathfrak c_T\mathfrak c_{ST}^{-1}
\]
are principal ideals \((a_{S,T})\) satisfying the conditions of Lemma 2. Hence there are principal ideals \((d_S)\) such that, after replacing
\[
  \mathfrak c_S\quad\text{by}\quad \mathfrak d_S=\mathfrak c_S/(d_S),
\]
their transformation systems become the unit ideal. Lemma 1 then gives an ideal \(\mathfrak b\) with
\[
  \mathfrak d_S=\mathfrak b^{1-S}.
\]
Passing to classes, the original classes \(c_S\) are also symbolic \((1-S)\)-st powers of the class of \(\mathfrak b\). This proves the theorem.

In the cyclic case the induced identity class becomes, after normalization, the group of principal ideals \((a)\) for which \(a\) is a local norm at all ramified places. Thus the theorem becomes the classical principal genus theorem.

\clearpage
\section{42. Split Crossed Products and Their Maximal Orders}

\textit{Original publication: Actualités scientifiques et industrielles 148 (1934), pp. 5--15.}

In the simplest case of split crossed products, that is, those whose factor systems are associated to one and hence give split algebras, the situation can be followed quite explicitly. The same remains true when the splitting field, or maximal commutative subfield, is not Galois. The elementary algebraic facts needed for this are developed first. I then give, for an algebraic number field as ground field, explicit descriptions of all maximal orders and of their ideals by means of modules and complementary modules of the splitting field. I also introduce a division into regions. Maximal orders are put into the same region when they have the same intersection with the chosen splitting field. The regions thus correspond bijectively to the orders in that field; the principal order corresponds to the principal region. The maximal orders in a region transform into one another by ideals built from modules of the corresponding order. In the principal region this is simply the extension of ideals from the splitting field, or of modules over its principal order.

Finally I pass to arbitrary simple algebras. By refining the division into regions - requiring not only the same intersection, but also equality of the local components at the ramified places - the theorems obtained in the split case transfer to the general case. In particular, the maximal orders of a principal region still transform into one another by extension ideals coming from the chosen maximal subfield.

\subsection{Matrix units of split crossed products and complementary bases}

Let \(\Omega\) be an arbitrary ground field. Let \(k/\Omega\) be a separable Galois extension of degree \(n\), with Galois group \(G\). Consider the crossed product with factor system one:
\[
  K=\sum_{S\in G} u_S k,
\]
with
\[
  z u_S=u_S z^S\qquad(z\in k),
\]
and
\[
  u_Su_T=u_{ST}.
\]
It is a split algebra. Put
\[
  E_1=\sum_{S\in G} u_S .
\]
Then \(E_1\) generates the identity representation of the group and satisfies the elementary relations
\[
  zE_1=E_1 z,
\]
for invariant \(z\), and more generally the trace relation
\[
  E_1 z E_1=E_1\operatorname{Tr}_{k/\Omega}(z).
\]
Here \(\operatorname{Tr}\) is the trace of the element \(z\) of \(k\) over \(\Omega\).

Let \(a_1,\ldots,a_n\) be a basis of \(k/\Omega\), and let \(a_1^*,\ldots,a_n^*\) be the complementary trace basis, so that
\[
  \operatorname{Tr}_{k/\Omega}(a_i a_j^*)=\delta_{ij}.
\]
Then
\[
  e_{ij}=a_i E_1 a_j^*
\]
form a complete system of matrix units in \(K\). Indeed,
\[
  e_{ij}e_{rs}=a_iE_1(a_j^*a_r)E_1a_s^*
  =a_iE_1\operatorname{Tr}(a_j^*a_r)a_s^*
  =\delta_{jr}e_{is}.
\]
Thus the split crossed product is explicitly the full matrix algebra, and the basis-complement relation is the device which supplies the matrix units.

The construction survives scalar extension. If \(Z\) is any commutative extension field of \(\Omega\), the same formulas hold after replacing \(k\) by \(kZ\). When \(kZ\) becomes a direct sum of fields, the group substitutions are defined to act trivially on \(Z\), in accord with the fact that the relevant elements commute elementwise with \(Z\). Thus the normal form is not a peculiarity of field extensions with one component; it is a statement about the separable algebra obtained by extension of scalars.

Assume now that \(\Omega\) has characteristic zero. The displayed normal form also holds when \(k\) is only a non-Galois splitting field. Let \(L\) be the Galois closure of \(k/\Omega\), with Galois group \(G\), and let \(H\) be the subgroup corresponding to \(k\). For the split crossed product built from \(L\) one forms the idempotent belonging to the identity representation of \(H\). Compressing by that idempotent gives back the algebra split by \(k\). The same trace-basis calculation then gives the matrix units. This proves, in particular, that a full matrix algebra over \(\Omega\) contains every degree-\(n\) field as a maximal subfield in a way compatible with the normal form.

The same passage from the Galois closure to the non-Galois splitting field works even when the algebra is not split. If \(L\) is a Galois splitting field and \(k\) corresponds to \(H\), the idempotent attached to \(H\) cuts out an algebra isomorphic to the original algebra with maximal subfield \(k\). It is the endomorphism ring of a suitable left ideal, and its rank is the right one because the left ideal decomposes into \(h=|H|\) mutually isomorphic summands after passing to the Galois closure. In the split case this reduces to the preceding construction.

\subsection{Maximal orders in the split case}

Let the ground field \(\Omega\) now be an algebraic number field. Let \(o\) be its maximal order, and let \(k/\Omega\) be a splitting field with maximal order \(\mathfrak o\). If \(\mathfrak a\) is a full \(o\)-module in \(k\), let \(\mathfrak a'\) denote the complementary module with respect to the trace form:
\[
  \mathfrak a'=\{x\in k: \operatorname{Tr}_{k/\Omega}(x\mathfrak a)\subseteq o\}.
\]
In the split algebra write again \(E_1\) for the idempotent belonging to the identity representation. The maximal orders and their ideals are described by the following formulas.

For a full module \(\mathfrak a\) put
\[
  \mathcal O_{\mathfrak a}=\mathfrak a E_1\mathfrak a'.
\]
Then \(\mathcal O_{\mathfrak a}\) is a maximal order. Conversely, every maximal order of the split algebra is obtained in this way. The left ideals between such maximal orders are exactly the modules
\[
  \mathfrak aE_1\mathfrak b,
\]
where \(\mathfrak a\) and \(\mathfrak b\) are full modules in \(k\). The reciprocal right ideal is
\[
  \mathfrak b'E_1\mathfrak a'.
\]
The intersection of \(\mathcal O_{\mathfrak a}\) with \(k\) is the coefficient order of \(\mathfrak a\). Hence the maximal orders whose intersection with \(k\) is a fixed order \(\mathfrak t\) form the region attached to \(\mathfrak t\).

The proof is local in nature. A module in the algebra is determined by its components at all places. If a module is contained in another at every completion, then it is contained globally; if it is properly enlarged, it is properly enlarged at some place. Hence maximality can be checked componentwise. At a local principal ideal ring choose complementary bases for \(\mathfrak a\) and \(\mathfrak a'\). The products
\[
  a_iE_1a_j^*
\]
are then matrix units, so the local order is the full matrix order and therefore maximal. The trace formula also gives
\[
  \operatorname{Tr}(\mathfrak a\mathfrak a')=o,
\]
which is the same assertion written in the language of complementary modules.

Let \(\mathcal O=\mathfrak aE_1\mathfrak a'\). A left ideal \(\mathcal L\) of \(\mathcal O\) is generated, as a module, by components of the form \(E_1\lambda\). Writing each \(\lambda\) in the splitting field gives
\[
  \mathcal L=\mathfrak aE_1\mathfrak b
\]
for some module \(\mathfrak b\). Regularity of the ideal is exactly the condition that \(\mathfrak b\) have full rank. Composition of such ideals corresponds to multiplication of the right-hand modules. In particular, two maximal orders in the same region transform into one another by ideals made from the modules belonging to the common coefficient order.

In the principal region, where the intersection with \(k\) is the maximal order \(\mathfrak o\), the statement becomes especially simple. The maximal orders in the principal region transform into one another by extension of ideals of \(k\). Equivalently, those left ideals of a maximal order in the principal region which are prime to the different and whose right order contains \(\mathfrak o\) are precisely the extended ideals from \(k\). This gives the arithmetic content of the explicit formulas above.

\subsection{Maximal orders of arbitrary crossed products}

The results just obtained remain true for an arbitrary simple algebra once the division into regions is refined. Let \(A\) be a central simple algebra over \(\Omega\), and let \(k\) be a maximal commutative subfield. Two maximal orders of \(A\) are assigned to the same region relative to \(k\) if they have the same intersection with \(k\) and, in addition, have the same local components at the ramified places of \(A\). If the common intersection is the maximal order of \(k\), the region is a principal region.

For principal regions one has the following theorem. The maximal orders of a principal region transform into one another by extension ideals of ideals from \(k\). In other words, the left ideals of maximal orders in a principal region which are prime to the different and whose right order contains the maximal order of \(k\) are precisely the extended ideals from \(k\).

The reason is that two maximal orders in a principal region can differ only at finitely many places, and by definition not at the ramified ones. At the remaining places the algebra is split. There it is of the form treated above, even if the splitting field is not Galois, by passage to the Galois closure and compression by the suitable idempotent. Thus the local ideals are extension ideals. Gluing the finitely many local data gives the global assertion.

If the algebra has prime degree, there is only one principal region. Indeed, locally at a ramified place it remains a division algebra, so no further ramified-place distinction can occur within the principal intersection condition. Hence the principal region is stable under extension ideals from \(k\). The most general transformation is obtained by composing these extension ideals with the ideals which transform a fixed maximal order into itself, namely the two-sided ideals. These differ from center ideals only at the ramified places; thus the statement about ideals of a fixed maximal order follows.

For arbitrary regions the corresponding assertion is local. The maximal orders of a region transform into one another by ideals prime to the different whose local components are built, at every place, from modules of the associated order in exactly the manner described in the split case. For a fixed ideal, the components differ from the order at only finitely many places. This is a direct consequence of the split calculation after replacing the operators \(u_S\) at the separate places by equivalent local ones.

There remains a finer question. Does every order in the maximal subfield occur as the intersection of a maximal order with that field? This is a question at the ramified places. Principal regions always exist, since a crossed-product presentation, or the generalized version of such a presentation, gives an order containing the maximal order of the subfield after the factor systems have been chosen integral.

One can also ask whether the arithmetic division into regions characterizes the algebraic splitting fields. Do distinct splitting fields always give distinct region divisions? Do they already give different principal regions? Do the principal regions belonging to all splitting fields exhaust all maximal orders? A single maximal order cannot provide such a characterization. Already in the quaternion algebra one finds examples where one maximal order contains both the principal order of a quadratic subfield and the principal order of the field of third roots of unity.

\clearpage
\section{43. Ideal Differentiation and the Different}

\textit{Original publication: J. reine angew. Math. 188 (1950), pp. 1--21. The paper was left by Noether and written in the winter of 1927/28; it was published posthumously.}

The theory of the different is usually developed in parallel with the differential calculus of algebraic equations. If an integral element \(x\) satisfies a defining equation \(f(x)=0\), the derivative \(f'(x)\) enters the different. This suggests that the different is not merely analogous to differentiation but is itself an ideal-theoretic form of differentiation. The following pages make this precise. The method begins with purely algebraic constructions for rings, then specializes to commutative rings and orders. The usual arithmetic different appears as the ideal derivative of a defining ideal.

The guiding idea is this. To extend the coefficient ring of a ring \(D\) from a base ring \(o\) to a larger ring \(h\), one has to construct a direct product of \(D\) with \(h\) over \(o\). If \(D\) is presented by generators and relations over \(o\), the same relations after coefficient extension define the product, provided the compatibility conditions are satisfied. The ideal which records these relations is the defining ideal. Differentiating this defining ideal, in the ordinary formal sense with respect to the generators, produces the different.

\subsection{Extension of the coefficient domain of a ring}

Let \(D\) and \(h\) be rings with common subring \(o\). A ring \(D_h\) is called the direct product of \(D\) with \(h\) over \(o\) if it is generated by isomorphic images of \(D\) and \(h\), if the two images have the same copy of \(o\), and if every element of \(D_h\) is a finite sum of products
\[
  x y\qquad(x\in D,\\ y\in h).
\]
The multiplication has to be bilinear over \(o\), and the two factors must commute with each other in the sense required by the construction. Thus, in the commutative case, one may think of
\[
  D_h=D\otimes_o h,
\]
but the definition is set up so as to apply before commutativity is imposed.

A necessary and sufficient condition for the direct product to exist is that all equality and multiplication relations among the chosen generators of \(D\) remain valid after the coefficients have been extended to \(h\), and similarly for the relations in \(h\). If \(D\) is given by generators \(Z_1,\ldots,Z_m\) over \(o\), let
\[
  o[Z_1,\ldots,Z_m]
\]
be the corresponding polynomial ring or non-commutative polynomial ring, as appropriate. The kernel of the natural map to \(D\) is the defining ideal \(M\). After coefficient extension the candidate direct product is
\[
  h[Z_1,
  \ldots,Z_m]/hM.
\]
It is the true direct product precisely when no new relation collapses the image of \(h\) or of \(D\) beyond those forced by \(M\).

Equivalently, the direct product has the following universal property. Given any ring \(R\) containing compatible images of \(D\) and \(h\) over \(o\), there is a unique homomorphism
\[
  D_h\longrightarrow R
\]
extending the two maps. In this form the construction is independent of the chosen generators. The defining ideal, however, is most convenient for calculation.

If \(D\) and \(h\) are commutative and \(D\) is finite over \(o\), the coefficient extension can be checked by bases. Let \(x_1,\ldots,x_n\) be elements of \(D\) which form an independent \(o\)-module basis. Then the products
\[
  x_ix_j=\sum_k c_{ij}^{k}x_k,
\]
with \(c_{ij}^k\in o\), determine the multiplication after extension to \(h\) by the same structure constants. This gives a direct product with \(h\)-basis \(x_1,
\ldots,x_n\), unless the extension of scalars introduces new linear dependence. Thus independence of the basis after extension is the essential condition.

For ideals and modules one uses the same language. If \(B\) is an \(o\)-module in \(D\), its extension is
\[
  B_h=B\otimes_o h
\]
inside the coefficient-extended ring. If \(C\) is an \(h\)-module in \(D_h\), its contraction is its intersection with \(D\). When \(B\) is an ideal, its extension is an ideal of \(D_h\); conversely, contraction of ideals is compatible with multiplication when the direct product is exact enough. This extension-contraction formalism is the module-theoretic counterpart of passing from an equation to its equations after scalar extension.

\subsection{Rings with an independent module basis}

Assume now that \(D\) has an independent basis \(x_1,
\ldots,x_n\) over \(o\). The direct product with \(h\) is constructed by taking all formal sums
\[
  \sum_i x_i a_i,
  \qquad a_i\in h,
\]
and multiplying them by the same structure constants as in \(D\). Associativity and distributivity follow from the corresponding identities in \(D\), because the basis is independent and the structure constants lie in \(o\). Thus the direct product exists and has \(h\)-basis \(x_1,
\ldots,x_n\).

Let \(B\) be an \(o\)-module in \(D\). Suppose \(B\) has an independent \(o\)-basis which can be enlarged to an independent basis of \(D\). Then extension and contraction behave without loss:
\[
  (B_h)\cap D=B.
\]
If \(B\) is an ideal, then \(B_h\) is an ideal, and contraction of \(B_h\) is again \(B\). This is the basic technical fact needed later. It is the reason that the different can be calculated after passing to a convenient coefficient extension.

A sufficient criterion for the existence of the direct product is therefore the following. If the ring and the modules under consideration possess independent bases over the coefficient ring, and if the multiplication tables have coefficients in that coefficient ring, then scalar extension gives the direct product. In arithmetic language this applies to orders in separable extensions and to the modules which occur as ideals of those orders.

\subsection{The different of a ring over a subring}

From here on the rings are assumed commutative unless explicitly stated otherwise. Let \(D\) be a commutative ring over \(o\), and suppose the direct product
\[
  D\otimes_o D
\]
exists in the sense described above. Multiplication in \(D\) gives a homomorphism
\[
  \mu:D\otimes_o D\longrightarrow D,
  \qquad x\otimes y\mapsto xy.
\]
The kernel of \(\mu\) is generated by the formal differences
\[
  x\otimes1-1\otimes x.
\]
The different is the ideal of \(D\) determined by the annihilator of this kernel. In modern notation this is the same as the zeroth Fitting ideal of the module of differentials, but Noether's construction is phrased entirely in terms of ideals and direct products.

Let \(D\) be given over \(o\) by generators \(Z_1,
\ldots,Z_m\) and defining ideal \(M\). For a polynomial \(F\in M\), form the formal differential
\[
  dF=\sum_i \frac{\partial F}{\partial Z_i}
  dZ_i.
\]
The ideal generated in \(D\) by the appropriate minors of the matrix of the partial derivatives of a system of generators of \(M\) is the ideal derivative of \(M\). This is the differential quotient of the defining ideal. The theorem is that this ideal derivative is exactly the different just defined by means of the direct product.

In the simplest case one generator is enough. Suppose
\[
  D=o[x],\qquad f(x)=0,
\]
and \(f(X)\) is a defining polynomial. Then the direct product is described by two elements \(x\) and \(x'\), with
\[
  f(x)=f(x')=0.
\]
The kernel of multiplication is generated by \(x'-x\). Since
\[
  f(x')-f(x)=(x'-x)\,q(x,x'),
\]
putting \(x'=x\) gives
\[
  q(x,x)=f'(x).
\]
Thus the different is generated by
\[
  f'(x),
\]
the ordinary derivative of the defining equation. The general case is the same calculation with several generators and the Jacobian matrix of the defining relations.

This proves that the classical derivative appearing in the different is not an accident of notation. It is the formal derivative of the defining ideal, transported through the direct-product description of coefficient extension.

\subsection{The different of a direct sum}

If a ring decomposes as a direct sum of components,
\[
  D=D_1\oplus\cdots\oplus D_r,
\]
then the different decomposes correspondingly. The direct product of \(D\) over \(o\) decomposes into the direct sum of the pairwise products \(D_i\otimes_o D_j\). The multiplication map is zero on the mixed components and is the ordinary multiplication on each diagonal component. Hence the kernel and its annihilator split by components. Consequently
\[
  \mathfrak D(D/o)=\mathfrak D(D_1/o)\oplus\cdots\oplus\mathfrak D(D_r/o),
\]
with the evident interpretation of the direct-sum ideal.

This observation is important because coefficient extension of a field often turns a field into a direct sum of conjugate fields. The different can then be computed componentwise. In particular, after a separable scalar extension to a splitting field, the structure of the direct product becomes diagonal, and the different is read off from the complementary bases of the components.

\subsection{Galois extension rings}

Let \(K/k\) be a finite Galois extension, and let \(D\) be an order in \(K\) over an order \(o\) of \(k\). After extending coefficients from \(o\) to \(D\), the direct product
\[
  D\otimes_o D
\]
decomposes into components corresponding to the automorphisms \(S\) of the Galois group. In the field case this is the familiar decomposition
\[
  K\otimes_k K \cong \prod_{S\in G} K,
\]
where the component indexed by \(S\) is obtained by the map
\[
  x\otimes y\longmapsto x\,y^S.
\]
For orders, one obtains modules inside the corresponding components.

The kernel of multiplication is the part which vanishes on the identity component. The other components measure the differences between an element and its conjugates. Thus the different is controlled by the product of the ideals generated by
\[
  x-x^S,
  \qquad S\ne 1,
\]
as \(x\) ranges through the order. This is the ideal-theoretic form of the discriminant calculation.

Equivalently, choose complementary bases \(\omega_1,
\ldots,
\omega_n\) and \(\omega_1^*,
\ldots,\omega_n^*\) for the trace pairing. Then the elements
\[
  \sum_i \omega_i\otimes \omega_i^{*S}
\]
are the component idempotents of the extended ring. The denominator ideals needed to make these idempotents integral are exactly the complementary module, and their inverse is the different. In this way the component representation of the Galois extension ring gives the standard formula
\[
  \mathfrak D^{-1}=\{x\in K: \operatorname{Tr}_{K/k}(xD)\subseteq o\}
\]
for the inverse different, and hence identifies the different itself with the inverse of the complementary module.

This is the structural core of the theory. The complementary module is not introduced as an auxiliary analytic object; it is the module of denominators needed to separate the components of the direct product. The different is its inverse.

\subsection{The different of an order}

Let \(D\) now be an order in a separable field extension \(K/k\), with coefficient order \(o\). The different \(\mathfrak D_{D/o}\) is the inverse of the complementary module
\[
  D^*=
  \{x\in K: \operatorname{Tr}_{K/k}(xD)\subseteq o\}.
\]
Thus
\[
  \mathfrak D_{D/o}=(D^*)^{-1}.
\]
If \(D\) has an independent \(o\)-basis \(\omega_1,
\ldots,
\omega_n\), the discriminant is
\[
  \Delta(\omega_1,
\ldots,
\omega_n)=
  \det(\operatorname{Tr}(\omega_i\omega_j))
\]
and the product of the basis ideal with the different gives the discriminant ideal. More precisely, changing the basis multiplies the determinant by the square of the determinant of the change-of-basis matrix, exactly as in the usual discriminant theory; the ideal formulation absorbs this dependence and gives an invariant ideal.

If the order is monogenic,
\[
  D=o[\theta],
\]
with fundamental equation \(f(\theta)=0\), then the different is generated by
\[
  f'(\theta)
\]
provided the basis \(1,\theta,
\ldots,
\theta^{n-1}\) is an integral basis of the order. In the non-monogenic case one replaces the one derivative by the ideal generated by the Jacobians of a system of defining relations. This is precisely the differential quotient of the defining ideal.

The formula is compatible with localization. At a prime ideal \(\mathfrak p\), localizing the order and the trace dual localizes the different:
\[
  (\mathfrak D_{D/o})_{\mathfrak p}
  =\mathfrak D_{D_{\mathfrak p}/o_{\mathfrak p}}.
\]
Therefore all assertions can be checked locally. Locally at a Dedekind prime, the different exponent is the valuation of the derivative in a primitive-element presentation, or equivalently the length of the module of differentials.

\subsection{Ramification theory of the principal order modulo powers of a prime}

The last part of the paper relates the ideal-theoretic different to ramification. Let \(\mathfrak p\) be a prime of \(o\), and let
\[
  \mathfrak pD=\prod_i \mathfrak P_i^{e_i}
\]
be its decomposition in the order or maximal order of the extension. The different exponent at \(\mathfrak P_i\) measures the failure of the residue extension and ramification to be unramified. In the tame case it is simply
\[
  e_i-1.
\]
In the wild case additional terms occur, represented by the higher ramification groups. Noether's formulation obtains these terms from the ideal derivative: the higher powers of \(\mathfrak P_i\) which divide all differences \(x-x^S\) are exactly the powers which enter the local different.

In a Galois extension the local different can therefore be written in terms of the ramification groups. If \(G_j\) denotes the usual decreasing filtration of the decomposition group, then the exponent of the different is
\[
  \sum_{j\ge0} (|G_j|-1),
\]
which is the Dedekind-Hensel formula. In the present framework this formula is not an external ramification theorem but the local expression of the same ideal-differentiation construction. The derivative of the defining ideal records, component by component in the direct product, the orders of contact of conjugate components; these orders of contact are exactly the ramification jumps.

Thus the different is simultaneously three objects: the annihilator attached to multiplication in the direct product, the ideal derivative of the defining relations, and the inverse of the trace-complementary module. The equality of these descriptions is the content of ideal differentiation.


\clearpage
\section{44. Algebra of Hypercomplex Quantities}

\textit{Lecture by Emmy Noether, winter semester 1929/30; written up by Max Deuring. This reader gives the lecture text through the applications of the cyclic special case.}

\subsection{Introduction}

These lectures treat the general representation theory of rings. This theory grew, on one side, out of the representation theory of finite groups and, on the other side, out of the algebraic and arithmetic theory of hypercomplex number systems. The more general point of view has an advantage over the older theory of Frobenius and Schur. One does not represent merely a group; one represents the group ring, and more generally a hypercomplex system or an arbitrary ring. Thus one may use the theory of rings and ideals. This frees representation theory from many unwieldy computations.

The use of general ring and ideal theory does more than simplify the subject. It also carries it further. It clarifies the relation of a group to its subgroups, and it opens new fields of application. In particular, general representation theory gives a new and elegant foundation of Galois theory; still more importantly, it gives a Galois theory for non-commutative fields.

The basic concepts of group, ring, and ideal theory, and the general theorems of representation theory, are those developed in Noether's paper \textit{Hypercomplex Quantities and Representation Theory}. The present lecture notes are to be read as a supplement to that work. The notions of representation and representation module are recalled briefly, chiefly in order to introduce reciprocal representation modules and reciprocal representations.

\subsection{Chapter I. Representations}

Let \(o\) be a ring, the ring to be represented, and let \(T\) be another ring, the ring in which the representation is to take place.

A direct representation of degree \(n\) of \(o\) in \(T\) is a subring \(O\) of the full matrix ring \(T_n\), together with a ring homomorphism from \(o\) onto \(O\). Symbolically,
\[
  o\twoheadrightarrow O\subset T_n.
\]
A reciprocal representation of degree \(n\) of \(o\) in \(T^*\) is defined similarly, but the map from \(o\) is a reciprocal ring homomorphism. Thus if \(c,d\in o\) correspond to \(c^*,d^*\), then \(c+d\) corresponds to \(c^*+d^*\), while \(cd\) corresponds to
\[
  d^*c^*.
\]
As an example, in a commutative coefficient field, transposition of matrices is a reciprocal ring isomorphism. In general, for every ring \(T\) one constructs a reciprocal ring \(T^*\) by assigning to each \(c\in T\) a symbol \(c^*\), putting
\[
  c^*+d^*=(c+d)^*,\qquad c^*d^*=(dc)^*.
\]
The symbols \(c^*\) then form a ring reciprocally isomorphic to \(T\).

Representations are collected into classes. A direct representation \(O\subset T_n\) is transformed into an equivalent representation by a regular matrix \(P\): instead of the matrix \(C\) attached to an element \(c\in o\) one uses
\[
  P^{-1}CP.
\]
All representations obtained from one another in this way form a direct representation class. The definition for reciprocal representations is the same.

A direct representation module of \(o\) with respect to \(T\) is an \(o,T\)-bimodule \(M\): it is a left \(o\)-module, a right \(T\)-module, and it satisfies the associativity law
\[
  (cm)t=c(mt)
\]
for \(c\in o\), \(m\in M\), and \(t\in T\). It is also required to be a finite direct sum of one-generator right \(T\)-modules,
\[
  M=x_1T+\cdots+x_nT,
\]
with the usual independence condition. A reciprocal representation module is a left \(o\)-module and a left \(T^*\)-module \(M^*\), satisfying
\[
  c(t^*m)=t^*(cm),
\]
and again decomposing as a finite direct sum of simple one-generator \(T^*\)-modules.

Every representation module determines a representation class. For a direct module \(M=x_1T+
\cdots+x_nT\), multiplication by an element \(c\in o\) is a \(T\)-linear transformation of \(M\) into itself, because
\[
  c(mt)=(cm)t.
\]
With respect to the basis \(x_1,
\ldots,x_n\) it is therefore given by a matrix \(C=(c_{ij})\), determined by
\[
  c x_j=\sum_i x_i c_{ij}.
\]
The matrices belonging to the elements of \(o\) form a homomorphic image of \(o\) in \(T_n\). The compatibility with addition and multiplication is immediate:
\[
  (c+d)x_j=cx_j+dx_j,
\]
\[
  (cd)x_j=c(dx_j).
\]
Thus a representation of degree \(n\) is obtained.

Changing the basis changes the matrices by similarity. If
\[
  (y_1,
\ldots,y_n)=(x_1,
\ldots,x_n)P
\]
with \(P\) regular, then the matrix \(C\) is replaced by
\[
  P^{-1}CP.
\]
Hence a module gives not a single representation but a whole representation class. Conversely, every representation class comes from a module. Starting with a representation \(c\mapsto C\), introduce a free right \(T\)-module
\[
  M=x_1T+
\cdots+x_nT
\]
and define left multiplication by \(o\) through
\[
  c(x_1,
\ldots,x_n)=(x_1,
\ldots,x_n)C.
\]
The homomorphism relations for the matrices are exactly the module axioms. Thus representation classes and representation modules are equivalent descriptions. Reciprocal representations and reciprocal modules behave in the same way, with the order of multiplication reversed.

\subsection{Chapter II. Galois theory of commutative fields}

The representation theory just developed is already sufficient to found Galois theory. Its advantage over the usual exposition is that it does not depend on special generating systems or on theorems about polynomials in indeterminates. Instead, it rests on coefficient extension of hypercomplex systems.

Let
\[
  A=c_1P+
\cdots+c_nP
\]
be a hypercomplex system over a commutative field \(P\), and let \(Z\) be an extension ring of \(P\) whose center contains \(P\). The extended system \(A_Z\) consists of the formal sums
\[
  \sum_i c_i z_i,
  \qquad z_i\in Z,
\]
with addition and scalar multiplication defined componentwise, and with multiplication determined by the multiplication table of \(A\). Thus if
\[
  c_ic_j=\sum_k a_{ij}^k c_k,
  \qquad a_{ij}^k\in P,
\]
then
\[
  \left(\sum_i c_i z_i\right)
  \left(\sum_j c_j w_j\right)
  =\sum_{i,j,k} c_k a_{ij}^k z_iw_j.
\]
This is the ordinary extension of coefficients.

Let now \(K/P\) be a commutative field extension of degree \(n\). If \(K\) is separable, then after a suitable extension of coefficients the algebra \(K\) decomposes as a direct sum of fields, one for each embedding. In representation-theoretic language, the regular representation decomposes into one-dimensional constituents. The automorphisms of the field are the isomorphisms of these constituents compatible with the multiplication.

For a field \(K\), an isomorphism over \(P\) is an operation preserving addition, multiplication, and the elements of \(P\). If \(K\) is of the first kind over \(P\), meaning that its scalar extension is completely reducible and has no radical, then the number of distinct \(P\)-isomorphisms into a normal closure equals the degree exactly when \(K/P\) is Galois. In that case the isomorphisms form a group.

The main theorem of Galois theory may then be stated without choosing a primitive element. Let \(K/P\) be a Galois extension with group \(G\). For a subfield \(T\) of \(K\), its invariant group is
\[
  G_T=\{S\in G: x^S=x\text{ for all }x\in T\}.
\]
For a subgroup \(H\subset G\), its invariant field is
\[
  K^H=\{x\in K: x^S=x\text{ for all }S\in H\}.
\]
Then:
\begin{enumerate}[label=\arabic*.]
\item If \(H\) is the invariant group of \(T\), then \(T\) is the invariant field of \(H\).
\item If \(T\) is the invariant field of \(H\), then \(H\) is the invariant group of \(T\).
\end{enumerate}
The first assertion follows from the extension theorem for isomorphisms: an isomorphism of a subfield can be extended to an isomorphism of the whole normal field. If an element of \(K\) is fixed by every automorphism fixing \(T\), then the one-dimensional component it defines in the extended regular representation is already determined by \(T\), and hence the element belongs to \(T\). The second assertion is reduced to the first by applying it to the invariant field of the subgroup.

This proof is deliberately structural. The usual arguments with equations are replaced by the decomposition of the extended system and by the behavior of isomorphisms under restriction and extension. It is precisely this form of the proof that can later be transferred to non-commutative fields.

\subsection{Chapter III. Abelian groups}

The next application concerns abelian groups and their group rings. Let \(G\) be a finite abelian group and let \(P\) be a field containing enough roots of unity, or else pass to an extension field which does. The group ring
\[
  P[G]
\]
is then a commutative semisimple algebra. Its irreducible representations are one-dimensional and are given by characters
\[
  \chi:G\longrightarrow P^{\times}.
\]
The algebra decomposes into a direct sum of fields corresponding to these characters.

The character relations are obtained from the matrix-unit or idempotent decomposition. For a character \(\chi\) define
\[
  e_\chi=\frac1{|G|}\sum_{g\in G}\chi(g^{-1})g.
\]
Then, when the characteristic does not divide \(|G|\), the \(e_\chi\) are primitive orthogonal idempotents and
\[
  1=\sum_\chi e_\chi.
\]
The orthogonality relations follow immediately:
\[
  \sum_{g\in G}\chi(g)\psi(g^{-1})=
  \begin{cases}
  |G|,&\chi=\psi,\\
  0,&\chi\ne\psi.
  \end{cases}
\]
No separate character calculation is needed; it is simply the formula for the decomposition of the group ring into primitive components.

The Galois theory of abelian groups is then the Galois theory of the commutative algebra \(P[G]\). Subgroups correspond to quotient character groups, and quotient groups correspond to subalgebras generated by the corresponding idempotent sums. In this form the duality between subgroups and character groups is just the duality between invariant rings and invariant groups already established.

If the field of coefficients is not large enough to contain the values of all characters, one extends coefficients to a splitting field and then descends. The Galois group of the coefficient extension permutes the primitive idempotents. The rational irreducible constituents over the original field are the sums over the orbits of this action. This is again the same method as before: extend to a completely reducible algebra, decompose, and then collect conjugate components.

\subsection{Chapter IV. Two-sided simple rings}

Let \(A\) be a two-sided simple hypercomplex system over a field \(P\). The basic lemma used throughout is the density statement for simple modules: if \(M\) is a simple representation module, then the ring of endomorphisms commuting with the action is a division ring, and the action of \(A\) on \(M\) is the full ring of endomorphisms over that division ring. Equivalently, a two-sided simple algebra is a full matrix algebra over a division algebra:
\[
  A\cong M_r(D).
\]
The center \(Z(D)\) is a field containing \(P\), and \(A\) is central simple over its center after replacing \(P\) by that center.

Representations of a two-sided simple system are controlled by its simple modules. Every representation module is a direct sum of isomorphic simple modules once the coefficient field is large enough to split the division algebra. Over the original field, the endomorphism division algebra measures the failure of the representation to be absolutely irreducible. The degree of this division algebra is the Schur index.

The passage from an arbitrary coefficient field to a splitting field has two effects. First, the center may decompose into conjugate components if it is not already a field after scalar extension. Second, each central simple component may become a full matrix algebra. The irreducible representations over the splitting field are then the natural column modules of these matrix algebras. Their conjugacy under the coefficient-field Galois group describes the irreducible representation classes over the original field.

This yields the non-commutative analogue of the field-theoretic Galois statements. In the commutative case a field is recovered as its own maximal commutative subfield. In the non-commutative case one must work with maximal commutative subfields and with their normalizers. The automorphism groups are replaced by groups of inner automorphisms, or by factor systems measuring the obstruction to representing automorphisms by elements satisfying the group law strictly.

\subsection{The Galois theory of non-commutative fields}

Let \(D\) be a division algebra over a field \(P\), and let \(H\) be a group of automorphisms of \(D\). The invariant field is
\[
  D^H=\{x\in D:x^S=x\text{ for all }S\in H\}.
\]
Conversely, for a subfield or subdivision ring \(E\subset D\), the invariant group consists of automorphisms fixing \(E\) elementwise. The aim is to obtain the same two assertions as in the commutative case:
\begin{enumerate}[label=\arabic*.]
\item the invariant group of \(E\) has invariant field \(E\);
\item the invariant field of \(H\) has invariant group \(H\), under the appropriate closedness condition.
\end{enumerate}
The proofs run parallel to the commutative proofs after replacing embeddings by simple components of representation modules and replacing ordinary automorphisms by inner automorphisms in a suitable simple algebra.

The normalizer of a maximal subfield plays the role of the Galois group. If \(K\) is a maximal commutative subfield of a central simple algebra \(A\), then an element \(u\in A^{\times}\) which normalizes \(K\),
\[
  u^{-1}Ku=K,
\]
induces an automorphism of \(K\). The obstruction to choosing such elements multiplicatively is measured by a factor system:
\[
  u_Su_T=u_{ST}a_{S,T},\qquad a_{S,T}\in K^{\times}.
\]
Thus non-commutative Galois theory naturally produces crossed products.

\subsection{Chapter V. Factor systems}

Let \(G\) be a group acting on a field \(K\). A factor system is a set of elements
\[
  a_{S,T}\in K^{\times}\qquad(S,T\in G)
\]
satisfying
\[
  a_{S,T}^{R}a_{ST,R}=a_{T,R}a_{S,TR}.
\]
This is exactly the associativity condition for symbols \(u_S\) with
\[
  u_Su_T=u_{ST}a_{S,T},
  \qquad z u_S=u_Sz^S.
\]
The resulting crossed product is
\[
  A=(a_{S,T},K,G)=\sum_{S\in G}u_SK.
\]

Changing the representatives \(u_S\) by
\[
  v_S=u_Sc_S
\]
changes the factor system to
\[
  a'_{S,T}=a_{S,T}c_S^Tc_Tc_{ST}^{-1}.
\]
Associated factor systems determine similar algebras. The classes of associated factor systems form an abelian group under componentwise multiplication. The identity consists of all transformation systems, i.e. all systems associated to the unit system. This group is the subgroup of the Brauer group consisting of classes split by \(K\).

A crossed representation with factor system \(a_{S,T}\) consists of matrices \(C_S\) over a coefficient ring satisfying
\[
  C_S^T C_T=C_{ST}a_{S,T}
\]
with the appropriate interpretation of the action of \(G\) on coefficients. Two such representations are equivalent if they differ by a change of basis. In this language ordinary representations are the special case \(a_{S,T}=1\). The representation theory of factor systems is therefore ordinary representation theory with multiplication twisted by a two-cocycle.

Multiplication of factor systems corresponds to tensor product of crossed representations. The inverse class is obtained by taking the reciprocal system. Powers of a factor-system class have a direct algebraic meaning: the \(m\)-th power corresponds to the tensor product of \(m\) copies of the associated crossed product. The index of the class is the degree of the division algebra in its Brauer class; it is also the degree of the irreducible crossed representations belonging to the system.

\subsection{Normal form with Galois maximal subfield}

Suppose a simple algebra \(A\) has a Galois maximal subfield \(K\) over its center \(P\). Then \(A\) can be written as a crossed product. For each automorphism \(S\) of \(K/P\) choose an element \(u_S\in A^{\times}\) inducing it by conjugation:
\[
  u_S^{-1}zu_S=z^S.
\]
Then every element of \(A\) has a unique expression
\[
  \sum_{S
\in G} u_S z_S,
  \qquad z_S\in K.
\]
The products of the \(u_S\) yield the factor system
\[
  u_Su_T=u_{ST}a_{S,T}.
\]
Associativity gives the cocycle relation. Conversely, any such crossed product is a central simple algebra with maximal subfield \(K\). Thus crossed products are not an artificial construction; they are the normal form of central simple algebras with Galois maximal subfield.

The normal form also clarifies splitting. The algebra is split if and only if the factor system is associated to one. Equivalently, there are elements \(c_S\in K^{\times}\) such that
\[
  a_{S,T}=c_S^Tc_Tc_{ST}^{-1}.
\]
Then the substitution \(v_S=u_Sc_S^{-1}\) makes
\[
  v_Sv_T=v_{ST},
\]
and the algebra is the endomorphism algebra of \(K\) over \(P\), hence a full matrix algebra.

\subsection{Chapter VI. Theory of crossed products}

Let \(K/P\) be Galois with group \(G\). The crossed product attached to a factor system \(a_{S,T}\) is the direct sum
\[
  A=\bigoplus_{S\in G}u_SK.
\]
It satisfies
\[
  z u_S=u_S z^S,
\]
\[
  u_Su_T=u_{ST}a_{S,T}.
\]
The center of \(A\) is \(P\), and \(K\) is a maximal commutative subfield. The algebra is simple. To see this, let \(I\) be a nonzero two-sided ideal and choose an element of \(I\) with the smallest possible number of nonzero components. Commuting it with elements of \(K\) eliminates all but one component unless the element already lies in \(K\). Since \(I\cap K\ne0\), and \(K\) is a field, \(1\in I\). Thus \(I=A\).

The product theorem for crossed products says that multiplying factor-system classes corresponds to taking the tensor product of the corresponding algebras over the common center and then passing to the component split by the diagonal copy of \(K\). In class notation,
\[
  [(a_{S,T})]+[(b_{S,T})]=[(a_{S,T}b_{S,T})].
\]
This gives the abelian group law on Brauer classes with fixed splitting field.

The automorphism ring of a suitable left ideal supplies the reciprocal algebra. If \(L\) is a left ideal of absolute length equal to the degree of \(K/P\), its endomorphism ring is a crossed product with inverse factor system. This gives the inverse in the Brauer group and explains why factor systems and their reciprocal systems cancel under multiplication.

If a factor system has index \(m\), its \(m\)-th power is split in the appropriate sense. Conversely the least positive exponent for which a power becomes split is tied to the exponent of the Brauer class, while the index is the degree of the underlying division algebra. In cyclic cases the two numbers can be read from norm groups; in general they are governed by the structure of the factor-system group.

\subsection{Applications of the cyclic special case}

The cyclic crossed product is especially explicit. Let \(K/P\) be cyclic of degree \(n\), generated by \(S\), and let
\[
  A=(a,K,S)
\]
be given by
\[
  A=K+uK+
\cdots+u^{n-1}K,
\]
\[
  zu=uz^S,
\]
\[
  u^n=a\in P^{\times}.
\]
Changing \(u\) to \(uc\) changes \(a\) to
\[
  aN_{K/P}(c).
\]
Thus the cyclic classes split by \(K\) are represented by
\[
  P^{\times}/N_{K/P}(K^{\times}).
\]

Several classical consequences follow. First, every finite field is commutative. Indeed, the multiplicative group of a finite field is cyclic; the center is a finite field; and the cyclic norm theorem forces the relevant crossed products to split. A finite division algebra over its center must therefore be commutative.

Second, finite Galois fields are cyclic over each of their subfields. The Frobenius automorphism generates the Galois group. In the crossed-product language this is the statement that the automorphism group controlling the extension is generated by one substitution, and all norm obstructions vanish because the multiplicative group is finite cyclic.

Third, for a \(p\)-adic ground field and a cyclic splitting field \(K/P\) of degree \(n\), the cyclic algebras are classified by the residue of \(a\) modulo the norm group. The local norm theorem determines exactly when
\[
  (a,K,S)
\]
splits. Hence every cyclic splitting field of degree \(n\) over a \(p\)-adic field gives the full set of division algebras whose invariant has denominator dividing \(n\). This is the local arithmetic form of the theory of crossed products.

The point of the lecture is that these results, which historically came from group characters, equations, and special computations, are all consequences of the same ring-theoretic mechanism: extend coefficients, decompose into simple components, identify the obstruction to strict multiplication by a factor system, and interpret the resulting class as an algebra class. In this form the theory is ready for the arithmetic applications to maximal orders, ideals, norm theorems, and the principal genus theorem.



\subsection{Further details from the lecture: two-sided simple systems}

We record the auxiliary lemma used for two-sided simple systems in a form suited to later applications. Let \(A\) be a simple system over \(P\), and let \(M\) be a simple left module. If elements \(m_1,
\ldots,m_r\) of \(M\) are linearly independent over the commuting division ring \(D=\operatorname{End}_A(M)\), and if \(n_1,
\ldots,n_r\) are arbitrary elements of \(M\), then there is an element \(a\in A\) with
\[
  a m_i=n_i\qquad(i=1,
\ldots,r).
\]
This is the familiar density theorem. In Noether's terminology it says that the representation system contains all linear substitutions allowed by the operator ring. The proof is the usual minimal-annihilator proof: one chooses a relation of least length among the conditions that fail, multiplies by suitable elements of \(A\), and contradicts the simplicity of the module.

From this lemma one obtains the structure theorem. If \(A\) is two-sided simple, then there is a division ring \(D\) and an integer \(r\) such that
\[
  A\cong D_{r}=M_r(D).
\]
The simple left modules are column modules of length \(r\) over \(D\), and the right operator ring of such a module is \(D\) itself. Conversely, the full matrix ring over a division ring is two-sided simple. This establishes the equivalence between simple systems and full matrix systems over skew fields.

Representations of a two-sided simple system are therefore governed by the decomposition of the coefficient extension of its center. Let \(Z\) be the center of \(A\). If \(Z/P\) is separable and \(\Omega\) is a splitting field, then
\[
  Z_\Omega=Z\otimes_P\Omega
\]
becomes a direct sum of conjugate fields. Consequently
\[
  A_\Omega=A\otimes_P\Omega
\]
becomes the direct sum of conjugate simple algebras, and over a sufficiently large \(\Omega\) each of these becomes a full matrix algebra. The irreducible constituents of the extended representation correspond to these components. If the center is inseparable, the extended center has a radical; passing to the quotient by the radical gives the same component decomposition for the semisimple quotient, and the radical contributes only composition factors operator-isomorphic to those already obtained.

Thus a representation irreducible over \(P\) is absolutely completely reducible after extension of scalars if and only if its center is separable over \(P\). If it is not separable, one must replace the direct-sum decomposition by a composition series; nevertheless the same conjugacy and multiplicity statements survive at the level of composition factors.

\subsection{Non-commutative fields and invariant theory}

A skew field \(D\) over a ground field \(P\) is treated as a simple system. Its subfields and automorphism groups are studied through their action on the associated representation modules. Let \(H\) be a group of automorphisms of \(D\). The invariant division ring is
\[
  D^H=\{x\in D: h(x)=x\text{ for every }h\in H\}.
\]
If \(E\subset D\) is a subdivision ring, its invariant group is
\[
  H_E=\{h\in\operatorname{Aut}(D):h(x)=x\text{ for every }x\in E\}.
\]
The analogy with ordinary Galois theory is exact only after adding a closedness condition. The invariant group of a subdivision ring is closed in the automorphism group. Conversely, for closed groups one recovers the group from its invariant ring.

The proof differs from the commutative proof only in language. One embeds the skew field in a full matrix algebra over a maximal commutative subfield and uses the theorem on extension of isomorphisms. The essential point is that any isomorphism of a suitable substructure can be extended to the whole simple system; in the matrix algebra this is an inner automorphism. Therefore the invariant operations are again described by normalizers. The formula
\[
  u^{-1}Eu=E
\]
for a normalizing element \(u\) replaces the equation \(S(E)=E\) of the commutative case.

If \(E\) is the invariant ring of \(H\), and \(u\) is an element which induces an automorphism fixing \(E\), then the extended-isomorphism theorem supplies an element of the group generated by \(H\) with the same action. Hence the automorphism lies in \(H\). This proves the non-commutative analogue of the Galois correspondence in the cases considered in the lecture.

\subsection{Factor systems and crossed representations in detail}

The lecture next develops factor systems as the exact algebraic measure of the failure of a choice of representatives to be multiplicative. Let \(K/P\) be Galois with group \(G\), and suppose that for every \(S\in G\) a representative \(u_S\) has been chosen. The representatives implement the automorphisms:
\[
  u_S^{-1} z u_S=z^S,
  \qquad z\in K.
\]
Since \(u_Su_T\) and \(u_{ST}\) implement the same automorphism, they differ by an element of \(K^{\times}\):
\[
  u_Su_T=u_{ST}a_{S,T}.
\]
Associativity of multiplication gives
\[
  (u_Su_T)u_R=u_S(u_Tu_R),
\]
hence
\[
  a_{S,T}^{R}a_{ST,R}=a_{T,R}a_{S,TR}.
\]
This is the factor-system relation.

The factor system depends on the choice of representatives. If
\[
  v_S=u_Sc_S,
  \qquad c_S\in K^{\times},
\]
then
\[
  v_Sv_T=v_{ST}a'_{S,T},
\]
where
\[
  a'_{S,T}=a_{S,T}c_S^Tc_Tc_{ST}^{-1}.
\]
The two systems are associated. Association is the natural equivalence relation, because it leaves the isomorphism class of the crossed product unchanged.

A crossed representation of \(G\) with factor system \(a\) is a family of transformations \(C_S\) satisfying the same multiplication law as the representatives:
\[
  C_S^T C_T=C_{ST}a_{S,T}.
\]
If \(a=1\), this is an ordinary representation. If \(a\) is a transformation system, then the crossed representation can be untwisted by a change of basis and again becomes ordinary. Thus the representation theory depends only on the associated factor-system class.

The irreducible crossed representations belonging to \(a\) are precisely the simple modules of the crossed product \((a,K,G)\). If the corresponding algebra is a matrix algebra of degree \(r\) over a division algebra \(D\), then the irreducible crossed representation has degree equal to the degree of \(D\). This is the Schur index. Hence the Schur index is not an extra numerical invariant attached to characters; it is the degree of the division algebra determined by the factor system.

\subsection{Multiplication of factor systems and the Brauer group}

Let \(a=(a_{S,T})\) and \(b=(b_{S,T})\) be factor systems for the same Galois extension \(K/P\). Their product
\[
  (ab)_{S,T}=a_{S,T}b_{S,T}
\]
is again a factor system. If \(A=(a,K,G)\) and \(B=(b,K,G)\), then the tensor product \(A\otimes_PB\) contains a component whose factor system is \(ab\). Passing to similarity classes gives
\[
  [A]+[B]=[(ab,K,G)].
\]
This is the group law.

The identity class is represented by the unit factor system. It is also represented by any transformation system
\[
  a_{S,T}=c_S^Tc_Tc_{ST}^{-1}.
\]
The inverse of \(a\) is represented by the reciprocal factor system, or equivalently by the opposite algebra. Thus the set of factor-system classes split by \(K\) becomes an abelian group, which embeds in the Brauer group of \(P\).

Noether's formulation emphasizes that all this is already contained in ideal theory. A factor system defines a group extension
\[
  1\longrightarrow K^{\times}\longrightarrow G^*\longrightarrow G\longrightarrow1,
\]
where \(G^*\) consists of the cosets \(u_SK^{\times}\). Multiplication in the extension is determined by the factor system. Changing representatives changes the cocycle but not the extension. Multiplication of factor systems corresponds to Baer product of such extensions, and to tensor product of the algebras.

In arithmetic applications the elements \(a_{S,T}\) are replaced by ideals or ideal classes. The same formulas remain meaningful, because ideal multiplication is abelian and the Galois group acts on ideals. This is the bridge to principal genera and to maximal orders.

\subsection{Normal representation and minimal splitting fields}

Assume that \(A\) is a central simple algebra with Galois maximal subfield \(K\). The crossed-product normal form gives a representation of \(A\) on the right \(K\)-module with basis \(u_S\). The action of \(K\) is diagonal through the conjugates, while the action of \(u_S\) permutes the components and multiplies them by factor-system elements. This is the normal representation.

If the factor system is associated to one, the normal representation decomposes into ordinary matrix units. Explicitly choose \(c_S\) with
\[
  a_{S,T}=c_S^Tc_Tc_{ST}^{-1}.
\]
Then \(v_S=u_Sc_S^{-1}\) satisfies
\[
  v_Sv_T=v_{ST}.
\]
Putting
\[
  E=\sum_{S\in G}v_S
\]
one obtains idempotents and matrix units from complementary bases exactly as in the split case. Hence the algebra is a full matrix algebra.

A field \(L\) is a splitting field for \(A\) if
\[
  A\otimes_PL
\]
is a full matrix algebra over \(L\). A detachment field is a field over which the irreducible representation decomposes into absolutely irreducible constituents. Minimal splitting and detachment fields are described by embedding maximal commutative subfields into the simple system. Their degrees are multiples of the product of the degree of the center and the index of the algebra.

If the center is separable, the absolutely irreducible constituents are permuted by the Galois group of the center. If the center is inseparable, one works with the semisimple quotient after extending scalars and accounts for the radical by composition factors. In either case the normal representation encodes the field-theoretic data through the algebra's center and index.

\subsection{Crossed products as algebras and ideals}

The crossed product \(A=(a,K,G)\) can be recovered from the group extension \(G^*\) by adjoining addition. Conversely, from \(A\) one recovers \(G^*\) as the set of regular elements normalizing \(K\):
\[
  G^*=\bigcup_{S\in G}u_SK^{\times}.
\]
The quotient of \(G^*\) by \(K^{\times}\) is \(G\). Thus the multiplicative structure of the algebra already contains the Galois group and the factor system.

Ideals in crossed products are treated by the same method. A left ideal \(L\) of the crossed product gives an endomorphism ring
\[
  \operatorname{End}_A(L),
\]
which is a similar algebra. If \(L\) has the minimal possible absolute length, its endomorphism ring is the division algebra opposite to the one represented by \(A\). This supplies the inverse Brauer class. Products of ideals correspond to multiplication of factor-system classes. The ideal-theoretic view is therefore not separate from the factor-system view; it is its arithmetic realization.

When a maximal order is present, one chooses integral representatives of the factor system. Ideals of the maximal order transform the order into equivalent maximal orders. Locally this is Brandt's theory; globally it becomes the theory used in the principal genus theorem. The reason the same formulas appear in both places is that a factor system is simply the multiplication table of a normalizer, and an ideal class is the arithmetic substitute for a scalar multiplier.

\subsection{Cyclic algebras and norm residues}

For a cyclic extension \(K/P\) of degree \(n\) with generator \(S\), the crossed product has the form
\[
  A=(a,K,S)=K+uK+\cdots+u^{n-1}K,
\]
\[
  zu=uz^S,
\]
\[
  u^n=a.
\]
Changing \(u\) to \(uc\) changes \(a\) by a norm:
\[
  a\mapsto aN(c).
\]
Thus the class depends on the residue class of \(a\) modulo the norm group.

If \(a\) is a norm, say \(a=N(c)^{-1}\), then \(v=uc\) satisfies \(v^n=1\) and the algebra splits. Conversely, if the algebra splits, the unit factor system is associated to the cyclic factor system, so \(a\) is a norm. This is the algebraic form of the norm theorem for cyclic algebras.

Over a finite field every element of the ground field is a norm from a finite extension. Therefore every cyclic algebra over a finite field splits. Since every finite division algebra has a maximal commutative subfield which is finite and cyclic over the center, it follows that the division algebra is equal to its center. This is Wedderburn's theorem that every finite skew field is commutative.

Over a local field the norm group has finite index and is described by valuation and residue data. The cyclic algebra \((a,K,S)\) is split exactly when \(a\) lies in the local norm group. If \(K/P\) is unramified, every unit is a norm and the invariant depends only on the valuation of \(a\). This is the local invariant of the cyclic algebra. It is the same local computation used in the proof of the principal genus theorem: once the local factor system has unit parameter in an unramified extension, the local algebra splits.

\subsection{Closing synthesis of the lecture}

The lecture has passed through the following chain of ideas. A representation is best understood as a module. A module is best understood through its operator ring. Simple operator rings are matrix rings over division rings. Extending coefficients decomposes these rings into simple components. When a Galois group acts on the components, the obstruction to choosing the component isomorphisms multiplicatively is a factor system. A factor system is the multiplication law of a crossed product. Crossed products are exactly the central simple algebras with a chosen Galois maximal subfield. Their arithmetic is the arithmetic of their orders and ideals.

This chain turns the older character-theoretic representation theory into ring theory. It also explains why the same formulas occur in Galois theory, class-field theory, the theory of maximal orders, and the arithmetic of hypercomplex systems. Each occurrence is a different expression of the same mechanism: inner automorphisms replace external automorphisms, factor systems measure their failure to multiply strictly, and ideals measure the arithmetic ambiguity in choosing the representatives.

\end{document}
